% draft0: 序論のC1--C3に対応させる。ここで新しい結果は出さない。
\section{結論}
\label{sec:conclusion}
\begin{itemize}
    \item \textbf{書くこと：} 正常時ARの継続予測とBFによる順位付け（C1），正式条件で確認した性能・構成要素の役割（C2），頑健性と校正の限界（C3）を各一段落で回収する．
    \item \textbf{注意：} 結果が支持する範囲でRQ1--RQ3へ明答する．「常に最良」「原因を因果的に証明」「校正済み障害検知」は結論にしない．
    \item \textbf{今後の課題：} 外部データでの検証，正常時の変動を考慮した校正，未知開始時刻への拡張を優先順に短く示す．
    \item \textbf{TODO：} 正式方式・比較表の確定後に記入し，最後に序論・概要の主張と照合する．
\end{itemize}
